\section{Results and Discussion}

\begin{figure*}[t]
    \centering
    \figmaybe[width=0.92\linewidth]{../generated/figures/figure1_framework.pdf}
    \caption{\textbf{Calibrated geometry-consistency framework.} H001 treats
    relation predictors as relation sources, standardizes prediction rows,
    preserves object-pair identity, joins explicit 3D geometry, estimates
    \pgeom{}, and reports recall and violation operating points.}
    \label{fig:framework}
\end{figure*}

\begin{figure*}[t]
    \centering
    \figmaybe[width=0.92\linewidth]{../generated/figures/figure2_tradeoff.pdf}
    \caption{\textbf{Recall-violation tradeoff.} Lower \vAt{100} and higher
    \rAt{100} are better. Open3DSG is a Docker-reproduced averaged-BLIP
    second-source variant under the covered H001 scope, not a broad SOTA
    comparison.}
    \label{fig:tradeoff}
\end{figure*}

\begin{table*}[t]
\centering
\small
\setlength{\tabcolsep}{4pt}
\caption{\textbf{Main VL-SAT held-out result.} Probabilistic re-ranking
preserves recall while reducing geometric violations; the rule-verified
condition is a zero-violation diagnostic.}
\label{tab:vlsat-main}
\begin{tabular}{lcccc}
\toprule
Condition & \rAt{50} & \rAt{100} & \vAt{50} & \vAt{100} \\
\midrule
\texttt{semantic\_only} & 0.9599 & 0.9894 & 0.0247 & 0.0469 \\
\texttt{probabilistic\_recalibrated} & 0.9642 & 0.9921 & 0.0234 & 0.0391 \\
\texttt{rule\_verified\_point\_subtype} & 0.9587 & 0.9890 & 0.0000 & 0.0000 \\
\texttt{family\_specific\_p\_geom\_valid} & 0.9619 & 0.9914 & 0.0204 & 0.0310 \\
\bottomrule
\end{tabular}
\end{table*}

The central question is whether calibrated geometry consistency reduces
physically inconsistent relation predictions without collapsing recall.
Figure~\ref{fig:tradeoff} visualizes the main recall-violation operating
points, and Table~\ref{tab:vlsat-main} gives the corresponding VL-SAT values.
The reproduced \texttt{semantic\_only} ranking reaches \rAt{50}/\rAt{100} of
0.9599/0.9894 with \vAt{50}/\vAt{100} of 0.0247/0.0469. The main
\texttt{probabilistic\_recalibrated} setting improves \rAt{50}/\rAt{100} to
0.9642/0.9921 and reduces \vAt{100} to 0.0391. This supports the intended
recall-first use case: calibrated geometry validity can reduce violations
while preserving, and slightly improving, exact-label recall.

The stricter \texttt{family\_specific\_p\_geom\_valid} setting reaches
\rAt{50}/\rAt{100} of 0.9619/0.9914 and \vAt{50}/\vAt{100} of 0.0204/0.0310.
This provides a stronger violation-first operating point. The
\texttt{rule\_verified\_point\_subtype} setting achieves zero measured
violations while keeping \rAt{50}/\rAt{100} at 0.9587/0.9890, but we treat it
as a diagnostic rather than the default method because hard filtering exposes
one end of the reliability-recall tradeoff.

\begin{table}[t]
\centering
\small
\setlength{\tabcolsep}{3pt}
\caption{\textbf{Nontriviality controls.} The effect is not explained by
geometry-only ranking, distance-only ranking, shuffled geometry, or wrong-pair
geometry.}
\label{tab:controls}
\begin{tabular}{lccc}
\toprule
Control & \rAt{100} & \vAt{100} & Reading \\
\midrule
\texttt{geom-only} & 0.5049 & 0.0701 & no semantics \\
\texttt{distance-only} & 0.5642 & 0.0993 & heuristic \\
\texttt{shuffled-geom} & 0.9788 & 0.0559 & no identity \\
\texttt{wrong-pair} & 0.9788 & 0.0581 & wrong identity \\
\bottomrule
\end{tabular}
\end{table}

Table~\ref{tab:controls} reports the nontriviality controls. Geometry-only
ranking falls to \rAt{50}/\rAt{100} of 0.2028/0.5049 and has higher violations
than the main calibrated setting. Distance-only ranking also underperforms,
reaching \rAt{50}/\rAt{100} of 0.3835/0.5642. Shuffled geometry and wrong-pair
geometry preserve more of the source signal but still degrade both recall and
violation behavior relative to calibrated identity-preserving joins. These
controls support the design requirement that geometry evidence must be joined
to the correct object pair and calibrated as a relation-level reliability
signal.

\begin{table}[t]
\centering
\small
\caption{\textbf{GT-based verifier evaluation.} \pgeom{} separates GT-positive
relations from deterministic counterfactual negatives.}
\label{tab:gt-verifier}
\begin{tabular}{lcc}
\toprule
Metric & Rows & Value \\
\midrule
GT-positive nonviolated & 2,545 & 0.9972 \\
GT-negative nonsatisfied & 2,545 & 0.9694 \\
\pgeom{} AUROC & 5,090 & 0.9779 \\
\pgeom{} AUPRC & 5,090 & 0.9737 \\
\pgeom{} Brier & 5,090 & 0.0538 \\
\bottomrule
\end{tabular}
\end{table}

Table~\ref{tab:gt-verifier} tests the geometry signal independently from
prediction re-ranking. GT-positive rows remain nonviolated at a rate of 0.9972,
while GT-derived counterfactual negatives are nonsatisfied at a rate of 0.9694.
The calibrated \pgeom{} score separates these groups with AUROC/AUPRC of
0.9779/0.9737 and Brier score 0.0538. This supports the use of \pgeom{} as a
reliability signal while still leaving room for residual calibration risk.

Audit and visual sanity checks are supporting evidence rather than the primary
metric result. The structured audit covers 250 rows and reaches quality-issue
precision of 0.8933, while the reduced 50-row visual spot-check reaches a
target-bucket quality-issue rate of 0.9333 with contradiction rate 0.0333. We
use these checks to support failure-mechanism interpretation and to detect
annotation or verifier issues, not as a large-scale independent human audit.

\begin{table}[t]
\centering
\small
\setlength{\tabcolsep}{3pt}
\caption{\textbf{Open3DSG second-source evidence.} Results use a
Docker-reproduced averaged-BLIP variant under the covered H001-family scope.}
\label{tab:open3dsg}
\begin{tabular}{lcccc}
\toprule
Condition & \rAt{50} & \rAt{100} & \vAt{50} & \vAt{100} \\
\midrule
\texttt{semantic\_only} & 0.3945 & 0.4963 & 0.1326 & 0.1195 \\
\texttt{prob.\ recal.} & 0.3843 & 0.5580 & 0.0575 & 0.0803 \\
\texttt{rule\_verified} & 0.4149 & 0.5238 & 0.0000 & 0.0000 \\
\texttt{family\_specific} & 0.4530 & 0.5984 & 0.0228 & 0.0311 \\
\bottomrule
\end{tabular}
\end{table}

Table~\ref{tab:open3dsg} reports Open3DSG as second-source evidence rather than
a broad open-vocabulary SOTA claim. Under the measured H001-family scope,
Open3DSG \texttt{semantic\_only} ranking reaches \rAt{50}/\rAt{100} of
0.3945/0.4963 with \vAt{50}/\vAt{100} of 0.1326/0.1195. The
\texttt{probabilistic\_recalibrated} setting reduces violations to
0.0575/0.0803 while shifting recall to 0.3843/0.5580. The family-specific
setting reaches \rAt{50}/\rAt{100} of 0.4530/0.5984 with
\vAt{50}/\vAt{100} of 0.0228/0.0311. These results support a cross-source
relation-reliability claim within measured H001 families, with the
averaged-BLIP, filtered-split, covered-scope, exact-label denominator, and
residual-calibration caveats kept explicit.

\begin{figure*}[t]
    \centering
    \figmaybe[width=0.92\linewidth]{../generated/figures/figure3_geometry_panels.pdf}
    \caption{\textbf{Geometry-backed qualitative failures.} Open3DSG examples
    show proximity, relative-vertical, and support/contact cases where semantic
    plausibility conflicts with object-pair geometry. The rightmost case
    illustrates residual calibration risk.}
    \label{fig:failures}
\end{figure*}

The qualitative failure analysis explains why the metric changes occur and
provides candidate source rows for Figure~\ref{fig:failures}. Open3DSG failure
rows produce 57,736 real analysis rows with zero validation errors, and the
inspected 36-case qualitative queue shows that 23 cases are demoted by
geometry-aware re-ranking. Figure~\ref{fig:failures} uses geometry-backed
panels from this queue to show family-structured failures: proximity failures
often involve distance contradictions, relative-vertical failures involve
vertical-order contradictions, and support/contact failures expose float-gap or
support-plane contradictions.

The same qualitative analysis also exposes a limitation. Ten of the 36 sampled
rule-violated cases still have $\pgeom > 0.9$. This means the calibrated score
should not be interpreted as a hard validity label. It also justifies why the
paper reports probabilistic, rule-verified, and family-specific variants
separately rather than collapsing them into a single geometry-corrected output.

